\section{Empirical strategy}\label{sec:strategy}

For each documentary-reviewed revision $e$ at treated institution $i(e)$ with event year $E_e$, we build a stack of its own. The stack holds the revising institution and every institution with no threshold revision anywhere in the target window $[E_e-4,E_e+5]$, observed in the same calendar years whenever AUTM data exist. A comparison institution can appear in more than one stack if it is a valid non-reviser for several events. That is by design. Each revision date gets its own set of contemporaneous comparisons, and no institution that revised inside an event's window is used as a control for it.

We estimate the following specification separately for each principal outcome $m$:
\begin{equation}
Y^{m}_{ist}
=
\sum_{k\neq-1}\delta^{m}_k D^k_{ist}
+
\sum_{k\neq-1}\theta^{m}_k D^k_{ist}S_{e(s)}
+
\gamma^{m}\log R_{it}
+
\alpha^{m}_{is}
+
\lambda^{m}_{stg(i)}
+
\varepsilon^{m}_{ist},
\label{eq:stacked}
\end{equation}
where $i$ indexes institutions, $s$ event-specific stacks, and $t$ calendar years. The index $m$ picks one of the five dependent variables: licenses and options, the filing margin, new patent applications, patents issued, or invention disclosures. Count outcomes enter as $\log(1+y)$, and the filing margin is defined in equation~\eqref{eq:filingmargin}. $D^k_{ist}$ equals one for the treated institution in stack $s$ when it is $k$ years from its revision, and zero otherwise. $S_{e(s)}$ equals $+1$ for an upward revision and $-1$ for a downward one. Event year $-1$ is the omitted reference year.

The institution-by-stack fixed effects $\alpha^{m}_{is}$ absorb level differences within each comparison set that do not change over time. The term $\lambda^{m}_{stg(i)}$ is a set of stack-by-calendar-year fixed effects interacted with the institution's group $g(i)$, defined by research size and public or private control. Research-size terciles are based on each institution's median research expenditure over 1991--2023. $R_{it}$ is research expenditure. Each stack therefore compares the revising institution with contemporaneous non-revisers, while letting calendar-year shocks differ across broad types of institution. Building a separate comparison set for each event follows the logic used to avoid contaminated comparisons when treatment timing is staggered \citep{GoodmanBacon2021,CallawaySantAnna2021,SunAbraham2021,CengizEtAl2019}.

For outcome $m$, the path $\delta^{m}_k+\theta^{m}_k$ traces the treated-versus-control difference around upward revisions, and $\delta^{m}_k-\theta^{m}_k$ traces it around downward revisions. The difference between the two is $2\theta^{m}_k$. We summarize the six post-revision years with
\begin{equation}
\Delta_m=\frac{1}{6}\sum_{k=0}^{5}2\theta^{m}_k.
\label{eq:delta}
\end{equation}
Table~\ref{tab:main} reports a separate $\Delta_m$ for each of the five outcomes; we never pool the outcomes into a single dependent variable.

$\Delta_m$ is an observational quantity. The upward-minus-downward contrast compares how outcomes change after revisions in each direction, relative to contemporaneous non-revisers, and it nets out changes common to all revision events. It does not remove changes that come with the \emph{direction} of a revision. Reading it causally would require upward and downward revisers to have followed similar paths relative to their controls had there been no direction-specific change. Nothing guarantees that here, because universities choose how to revise. In particular, upward revisions start from predecessor documents with much lower PCSI than downward revisions, so mean reversion, differences between institutions, and different reasons for revising all remain serious concerns. We treat pre-revision paths and balance checks as evidence about comparability, not as proof that direction is exogenous.

Standard errors are clustered by institution. The finite-sample correction counts the slope coefficients but not the absorbed fixed effects; counting the stack-by-year-by-group effects as well, as some software does by default, raises standard errors by about 3--4 percent and changes no conclusion. Because one institution can appear in several stacks, the regression sample size counts stack-by-institution-by-year rows, not unique institution-years. The variation in revision direction comes from 34 events, only eight of them upward, not from the roughly 26,000 stacked rows. We therefore add direction-label permutation inference to the clustered standard errors. Holding event dates, stacks, outcomes, and controls fixed, we repeatedly choose which eight of the 34 events to label upward and recompute $\Delta_m$. The preferred results use 100,000 random assignments that keep the number of upward events fixed. The test asks whether the observed contrast is unusual if the direction labels are exchangeable across events. Direction was not randomly assigned, so that assumption has real content; the design does not guarantee it.

For each outcome we also report a joint test of the three pre-revision direction coefficients, leave-one-event-out estimates for licensing, and results for the stricter 29-event sample. Benjamini--Hochberg $q$-values adjust the five outcome-level permutation tests as one family. The adjustment matters here because licensing is the strongest of five outcomes we examine, not the only one.

The annual panel is supporting evidence, not the main design. We regress each outcome on contemporaneous and lagged PCSI with institution and year fixed effects, controlling for lagged log research expenditure, log licensing FTEs, and inventor royalty share, and clustering by institution. Most annual PCSI values are carried forward from older documents, and revision timing is blurred in that representation. The annual regressions therefore answer a broader question about association, and we read them separately from the documentary event study.
